\documentclass[11pt,a4paper]{article}

\usepackage[margin=1in]{geometry}
\usepackage{amsmath,amssymb,amsthm,mathtools}
\usepackage{hyperref}
\usepackage{booktabs}
\usepackage{microtype}
\usepackage{url}

\hypersetup{
    colorlinks=true,
    linkcolor=blue,
    urlcolor=blue,
    citecolor=blue
}

\title{A Digit-Avoidance Conjecture for Powers of Two}
\author{}
\date{October 2026}

\theoremstyle{definition}
\newtheorem{definition}{Definition}[section]

\theoremstyle{plain}
\newtheorem{theorem}[definition]{Theorem}
\newtheorem{lemma}[definition]{Lemma}
\newtheorem{proposition}[definition]{Proposition}
\newtheorem{corollary}[definition]{Corollary}

\theoremstyle{remark}
\newtheorem{remark}[definition]{Remark}
\newtheorem{conjecture}[definition]{Conjecture}
\newtheorem{question}[definition]{Open Question}

\newcommand{\forbidden}{\mathcal{D}}
\newcommand{\period}{P_k}

\begin{document}

\maketitle

\begin{abstract}
We investigate a digit-avoidance property of the powers of two
in base ten. The number
\[
2^{16}=65536
\]
contains none of the decimal digits \(1,2,4,8\). We conjecture
that it is the only nonnegative integral power of two with this
property.

Elementary modular arguments impose necessary restrictions on
any possible solution. We establish a periodicity result for
decimal suffixes and describe a computational method for
enumerating exponent residue classes that satisfy increasingly
long suffix restrictions. Two complementary Python scripts
provide a framework for searching individual powers and studying
the associated residue classes.

The uniqueness assertion remains open. Computational evidence
over finite ranges and patterns in residue-class counts are not
presented as proofs of the conjecture.
\end{abstract}

\section{Introduction}

The decimal expansion of an integer can exhibit properties that
are not immediately apparent from its algebraic representation.
For example,
\[
2^{16}=65536
\]
is a power of two whose decimal digits belong entirely to
\[
\{0,3,5,6,7,9\}.
\]
In particular, none of the digits \(1,2,4,8\) occurs.

This observation motivates a simple question: is \(65536\) the
only power of two whose decimal representation avoids all four
of these digits?

Although the question is elementary to state, proving uniqueness
requires controlling decimal digits over arbitrarily many positions.
The elementary results below restrict the possible exponents.
We then establish a modular framework for studying longer decimal
suffixes and describe computational experiments motivated by it.

\section{Definitions and Main Conjecture}

Let
\[
\forbidden=\{1,2,4,8\}.
\]

\begin{definition}
For \(n\in\mathbb{Z}_{\geq 0}\), let \(P(n)\) denote the statement
that the standard decimal representation of \(2^n\) contains no
digit belonging to \(\forbidden\).
\end{definition}

The conjectured example is
\[
2^{16}=65536,
\]
so \(P(16)\) holds.

\begin{conjecture}[Uniqueness conjecture]
For every \(n\in\mathbb{Z}_{\geq 0}\),
\[
P(n)\iff n=16.
\]
\end{conjecture}

The case \(n=0\) is not a solution, since \(2^0=1\). The case
\(n=16\) establishes existence. The remaining task is to prove
that no other exponent satisfies \(P(n)\).

\section{Restrictions from the Final Decimal Digit}

For \(n\geq 1\), the final decimal digit of \(2^n\) follows the
periodic cycle
\[
2,\ 4,\ 8,\ 6.
\]
Equivalently,
\[
2^n\bmod 10 =
\begin{cases}
2,&n\equiv 1\pmod 4,\\
4,&n\equiv 2\pmod 4,\\
8,&n\equiv 3\pmod 4,\\
6,&n\equiv 0\pmod 4.
\end{cases}
\]

Since \(2,4,8\) are forbidden, any positive solution must end
in the digit \(6\).

\begin{proposition}
If \(n\geq 1\) and \(P(n)\) holds, then
\[
n\equiv 0\pmod 4.
\]
\end{proposition}

\begin{proof}
The final digit of \(2^n\) must be \(6\). The final-digit cycle
shows that this occurs precisely when \(n\equiv 0\pmod 4\).
\end{proof}

Consequently, every positive candidate exponent is a multiple
of four. Writing \(n=4m\), we obtain
\[
2^n=2^{4m}=16^m.
\]

Thus the uniqueness conjecture is equivalent to the assertion
that \(16^m\) contains a forbidden decimal digit for every positive
integer \(m\neq 4\).

\section{Restrictions from the Final Two Decimal Digits}

The powers of \(16\) modulo \(100\) satisfy
\[
\begin{aligned}
16^1&\equiv16\pmod{100},\\
16^2&\equiv56\pmod{100},\\
16^3&\equiv96\pmod{100},\\
16^4&\equiv36\pmod{100},\\
16^5&\equiv76\pmod{100},\\
16^6&\equiv16\pmod{100}.
\end{aligned}
\]

Hence the final two digits cycle through
\[
16,\ 56,\ 96,\ 36,\ 76
\]
with period five.

When \(m\equiv 1\pmod 5\), the final two digits of \(16^m\)
are \(16\), which contains the forbidden digit \(1\). Therefore,
a solution must satisfy \(m\not\equiv 1\pmod 5\).

\begin{proposition}
If \(n\geq 1\) and \(P(n)\) holds, then
\[
n\equiv 0,8,12,\text{ or }16\pmod{20}.
\]
\end{proposition}

\begin{proof}
The preceding proposition gives \(n=4m\). The final-two-digit
cycle excludes \(m\equiv 1\pmod 5\), leaving
\[
m\equiv 0,2,3,\text{ or }4\pmod 5.
\]
Multiplying these congruences by four yields the result.
\end{proof}

\begin{remark}
These restrictions are necessary, not sufficient. Avoiding
forbidden digits in the final two positions does not imply that
the other digits are also allowed.
\end{remark}

\section{Periodicity of Decimal Suffixes}

We now consider the last \(k\) decimal digits of \(2^n\).
For a positive integer \(k\), define
\[
P_k=4\cdot 5^{k-1}.
\]

We first establish the relevant multiplicative-order formula.

\begin{lemma}
For every integer \(k\geq 1\),
\[
\operatorname{ord}_{5^k}(2)=4\cdot 5^{k-1}.
\]
\end{lemma}

\begin{proof}
Since \(2\) has order \(4\) modulo \(5\), any exponent \(r\)
satisfying \(2^r\equiv 1\pmod{5^k}\) must be divisible by
four. Write \(r=4m\).

For \(m\geq 1\), the standard lifting-the-exponent formula gives
\[
v_5(16^m-1)
=
v_5(16-1)+v_5(m)
=
1+v_5(m),
\]
where \(v_5(x)\) denotes the exponent of \(5\) in the prime
factorization of \(x\).

Since \(2^{4m}=16^m\), the congruence
\[
2^{4m}\equiv1\pmod{5^k}
\]
holds precisely when \(1+v_5(m)\geq k\), or equivalently
when \(5^{k-1}\mid m\).

The smallest positive such exponent is therefore
\(4\cdot5^{k-1}\), as required.
\end{proof}

\begin{theorem}[Periodicity of the last \(k\) digits]
For every \(k\geq1\) and every \(n\geq k\),
\[
2^{n+P_k}\equiv 2^n\pmod{10^k},
\qquad P_k=4\cdot5^{k-1}.
\]
Consequently, the last \(k\) decimal digits of \(2^n\) depend
only on \(n\bmod P_k\), once \(n\geq k\).
\end{theorem}

\begin{proof}
Because \(n\geq k\),
\[
2^n\equiv 0\pmod{2^k}.
\]
The same is true of \(2^{n+P_k}\). Thus
\[
2^{n+P_k}\equiv 2^n\pmod{2^k}.
\]

By the lemma,
\[
2^{P_k}\equiv1\pmod{5^k}.
\]
Multiplying by \(2^n\) gives
\[
2^{n+P_k}\equiv2^n\pmod{5^k}.
\]

Since \(2^k\) and \(5^k\) are coprime, the Chinese remainder
theorem yields
\[
2^{n+P_k}\equiv2^n\pmod{10^k}.
\]
Repeated application establishes the claimed dependence on
the exponent modulo \(P_k\).
\end{proof}

\begin{remark}
The integer \(2\) is not invertible modulo \(10^k\). The
periodicity result therefore does not follow by treating \(2\)
as an element of the multiplicative group modulo \(10^k\).
Instead, the proof uses the vanishing of \(2^n\) modulo \(2^k\)
for \(n\geq k\), together with the multiplicative order of
\(2\) modulo \(5^k\).
\end{remark}

\begin{remark}
The theorem establishes a period; it does not assert that
\(P_k\) is the smallest possible period of the suffix sequence.
When interpreting suffixes shorter than the full decimal
representation, leading zeros may be used to pad the suffix
to length \(k\). Such padding introduces no forbidden digits.
\end{remark}

\section{Admissible Residue Classes}

The periodicity theorem allows the suffix restriction to be
studied through a finite set of exponent residue classes.

\begin{definition}
For \(k\geq1\), let \(S_k\) be the set of residue classes
modulo \(P_k=4\cdot5^{k-1}\) whose representatives \(n\geq k\)
have last \(k\) decimal digits containing no digit from
\(\forbidden\). The periodicity theorem ensures that this
property is independent of the choice of representative,
provided the representative is at least \(k\).
\end{definition}

Write
\[
s_k=\lvert S_k\rvert.
\]
There are \(P_k\) residue classes in total, so
\[
0\leq s_k\leq P_k.
\]

Any exponent satisfying the original digit-avoidance property
must belong to an admissible suffix class for every suffix length
\(k\leq n\). However, admissibility for a fixed suffix length
does not establish that the entire decimal expansion is admissible.

The sequence of class counts \(s_k\) provides a way to investigate
how restrictive the suffix conditions become as \(k\) increases.

\section{A Conjectured Residue-Class Counting Pattern}

Computational experiments using \texttt{trailing\_classes.py}
are designed to count the admissible classes as \(k\) increases.
A pattern suggested by the preliminary analysis is
\[
s_1=1,\qquad
s_k=4\cdot3^{k-2}\quad(k\geq2).
\]

This formula is recorded as a conjecture about the class counts,
not as a theorem. A complete proof would require a rigorous
description of how admissible classes at level \(k\) lift to
classes at level \(k+1\).

In particular, the fact that each class has five possible
extensions does not by itself imply that exactly three extensions
remain admissible. The new decimal digit depends on the extension,
and the number of allowed digits depends on whether the digit
is odd or even. Any proposed recurrence must account for this
structure.

The script provides a finite computational test of the proposed
pattern for the values of \(k\) it examines. Its output should be
reported alongside the precise run parameters when presenting
numerical evidence. Such evidence does not substitute for a proof
of the formula for all \(k\).

Even if the proposed formula were established, it would describe
admissible suffix classes rather than directly count powers of two
whose entire decimal expansions avoid the forbidden digits.
The distinction is essential: a number can pass a suffix test
and still contain a forbidden digit farther to the left.

\section{Computational Methods}

Two complementary Python scripts accompany this investigation.
The source files are available in the repository:

\begin{itemize}
    \item \href{https://github.com/IronMan1405/mathematical_proofs/blob/main/scripts/pow2_digits.py}
    {\texttt{pow2\_digits.py}: exact search and suffix sieve}.
    \item \href{https://github.com/IronMan1405/mathematical_proofs/blob/main/scripts/trailing_classes.py}
    {\texttt{trailing\_classes.py}: residue-class enumeration}.
\end{itemize}

\subsection{Exact search with suffix filtering}

The script \texttt{pow2\_digits.py} implements two methods.

\begin{enumerate}
    \item \textbf{Direct search.} For each exponent in a specified
    finite interval, the program constructs \(2^n\) exactly and
    checks its full decimal representation.

    \item \textbf{Suffix sieve followed by exact verification.}
    The program tracks \(2^n\bmod 10^K\) incrementally. If the
    resulting \(K\)-digit suffix contains a forbidden digit, the
    exponent is rejected. Every candidate that passes the suffix
    test is then checked against the full integer \(2^n\).
\end{enumerate}

The sieve is logically exact. If a forbidden digit occurs among
the last \(K\) digits, the full decimal representation cannot
satisfy the conjectured property. Passing the sieve, however,
is only a necessary condition, so the full check remains essential.

The program compares the direct and sieved methods on a common
range and asserts that their reported solutions agree there.
This cross-check helps detect implementation errors, although
it does not prove anything about exponents outside the searched
range.

\subsection{Residue-class enumeration}

The script \texttt{trailing\_classes.py} enumerates exponent
residue classes modulo \(P_k\) and retains classes whose last
\(k\) digits avoid the forbidden set. It extends the classification
from one suffix length to the next by examining the possible
extensions of each surviving class.

This computation investigates the structure of the suffix
restriction rather than checking only individual powers. The
number of surviving classes at each level can be compared with
the proposed counting pattern in the preceding section.

Both scripts are finite computations. Their conclusions must
be limited to the exponent ranges and suffix lengths actually
examined.

\subsection{Reproducibility}

The scripts can be run with Python 3 using commands of the form
\begin{verbatim}
python pow2_digits.py 20000 2000000
python trailing_classes.py 13
\end{verbatim}

The first command sets the direct-search limit to \(20{,}000\)
and the sieved-search limit to \(2{,}000{,}000\). The second
requests residue-class counts through suffix length \(13\).
These are example commands, not a claim that a particular
execution has been performed or independently verified here.

For a reproducible report, the repository should preserve the
script versions, the exact command-line arguments, the Python
version, and the resulting output. Any numerical claims in the
paper should match that recorded output.

\begin{remark}
A finite search can establish
\[
P(n)\implies n=16
\qquad\text{for }0\leq n\leq N
\]
for a specified bound \(N\), provided every exponent in that
range is checked correctly. It cannot establish the same
statement for all nonnegative integers \(n\).
\end{remark}

\section{Further Questions}

The elementary restrictions leave infinitely many candidate
exponents. A proof of uniqueness would need to control the
decimal digits across arbitrarily many positions or introduce
another argument that excludes every exponent other than \(16\).

The modular approach suggests several questions:

\begin{enumerate}
    \item Can the admissible suffix classes be described by an
    exact recurrence as the suffix length increases?

    \item Is the observed formula
    \(s_k=4\cdot3^{k-2}\) valid for every \(k\geq2\), and if so,
    can it be proved?

    \item Can bounds on the number of admissible suffix classes
    yield meaningful restrictions on full powers of two?

    \item Can any such counting argument establish uniqueness,
    or would it only show that admissible exponents are sparse?

    \item Is there a useful connection with established results
    on the distribution of digits in sequences such as \(2^n\)?
\end{enumerate}

A density estimate, even if proved, would not alone establish
uniqueness. A set can have density zero and still contain
infinitely many elements. The central challenge remains to
exclude every exponent other than \(16\).

\section{Conclusion}

The identity
\[
2^{16}=65536
\]
provides an example of a power of two whose decimal expansion
avoids the digits \(1,2,4,8\).

Elementary modular arguments show that any positive solution
must have exponent divisible by four. An analysis of the final
two digits further restricts the exponent to the residue classes
\[
n\equiv0,8,12,\text{ or }16\pmod{20}.
\]

The periodicity theorem provides a rigorous foundation for
studying longer decimal suffixes through finite residue-class
enumerations. The accompanying scripts offer reproducible
methods for exploring both individual powers and suffix-class
counts.

Nevertheless, neither finite computation nor the observed
counting pattern establishes uniqueness over all exponents.
The central conjecture remains open:

\[
\boxed{
\text{\(65536\) is the unique nonnegative power of two
avoiding the decimal digits \(1,2,4,8\).}
}
\]

\end{document}